%% arXiv preprint starter — the widely used NeurIPS-derived preprint style.
%%
%% arxiv.sty IS bundled with this download: it is a standalone styling
%% package (not a CTAN package), distributed from its own GitHub repo, so
%% there is no TeX Live copy to fall back to.
%%
%% Source: https://github.com/kourgeorge/arxiv-style
%% Licence: MIT, (c) George Kour -- verified in the repo's own License.txt.
%%
%% This is a \usepackage, not a document class: it restyles the plain
%% article class rather than replacing it.
\documentclass{article}

\usepackage[T1]{fontenc}
\usepackage[utf8]{inputenc}
\usepackage{arxiv}
\usepackage{hyperref}
\usepackage{booktabs}
\usepackage{amsfonts}
\usepackage{microtype}
\usepackage{graphicx}
\usepackage{natbib}
\usepackage{doi}

\title{Your Title Here}

\author{
  Ada Lovelace \\
  Institute of Analytical Engines\\
  London, United Kingdom \\
  \texttt{ada@example.edu} \\
  \And
  Second Author \\
  Another University\\
  Springfield, USA \\
  \texttt{second@example.org} \\
}

\begin{document}
\maketitle

\begin{abstract}
Replace this with your abstract. A good abstract states the problem, what
you did, and what you found --- in that order, and without citations.
\end{abstract}

\keywords{First keyword \and Second keyword \and Third keyword}

\section{Introduction}
Start here. Cite things like this~\citep{knuth1984texbook}.

\section{Method}
\label{sec:method}
Describe what you did.

\begin{table}[htbp]
  \caption{A caption goes above the table.}
  \label{tab:example}
  \centering
  \begin{tabular}{lrr}
    \toprule
    Condition & Count & Rate \\
    \midrule
    Baseline  & 120   & 0.42 \\
    Ours      & 118   & 0.71 \\
    \bottomrule
  \end{tabular}
\end{table}

\section{Results}
Table~\ref{tab:example} shows the numbers.

\section{Conclusion}
Say what it means.

\bibliographystyle{unsrtnat}
\bibliography{references}

\end{document}
